\section{候选组成与排序边界分析}
\label{sec:analysis}

来源对照同时改变了候选身份与频率组成。本节依次检验总体差异是否掩盖分层变化，以及增加训练候选数量是否持续改善排序；两项分析分别界定收益分布与候选预算边界。

\subsection{总体接近是否掩盖频率层级差异}
\label{sec:frequency-slices}

表~\ref{tab:fb-paired}并列报告纠错和退化。Full与配额匹配组共同正确8,642题、共同错误669题，另有348题仅Full正确、341题仅配额匹配组正确。总体EM接近仍伴随689题判定变化，需要进一步检查其分布。

\begin{table}[htbp]
\caption{FB同10,000题上的配对变化。纠错表示左侧方法正确、右侧错误，退化为相反情况；净变化为纠错减退化。}
\label{tab:fb-paired}
\centering\small
\begin{tabular}{lrrr}
\toprule
比较方向 & 纠错 & 退化 & 净变化 \\
\midrule
Full $-$ 均匀随机 & 553 & 211 & +342 \\
Full $-$ Profile-only & 755 & 254 & +501 \\
Full $-$ Violation-only & 582 & 301 & +281 \\
Full $-$ 配额匹配随机 & 348 & 341 & +7 \\
配额匹配 $-$ 均匀随机 & 534 & 199 & +335 \\
\bottomrule
\end{tabular}
\end{table}


为检查总体净差是否掩盖不同频率关系上的相反变化，我们按训练关系词表中各关系边数的三分位次序统计量定义低、中、高频组，同频次不拆分；两条边界为215和689。边界仅由训练图决定，不按测试正确率调整。测试问题按目标关系的训练频次归组，所有10,000题均纳入，覆盖215个测试中出现的训练关系。该分析是在总体结果观察后的描述性分析，不构成独立确认性检验。

\begin{table}[htbp]
\caption{按目标关系训练边数划分的FB描述性配对分析。边界为训练关系词表频次的三分位次序统计量，同频次不拆分；未按测试表现调边界。EM单位为\%；配额匹配指Full配额匹配随机。所有分组均报告，净变化合计为7题，不构成跨种子或因果证据。}
\label{tab:fb-frequency-slices}
\centering\small
\begin{tabular}{lrrrrrr}
\toprule
目标关系频率 & $n$ & Full EM & 配额匹配EM & 纠错 & 退化 & 净变化 \\
\midrule
低频（$\leq215$） & 412 & 60.92 & 65.05 & 24 & 41 & -17 \\
中频（216--689） & 1089 & 82.74 & 82.64 & 51 & 50 & +1 \\
高频（$>689$） & 8499 & 92.22 & 91.95 & 273 & 250 & +23 \\
\bottomrule
\end{tabular}
\end{table}


低频组的412题中，Full相对配额匹配组纠错24题、退化41题，净减少17题（$-4.13$个百分点）；中频组净增加1题，高频组净增加23题。因此，当前数据不支持Full在低频关系上的补偿性优势。高频组占测试题目的84.99\%，按题等权的总体EM与按关系等权的结果可能不同。以测试中出现的关系等权平均EM，Full为75.21\%，配额匹配组为76.15\%；各方法的完整分层及关系平均结果见附录表~\ref{tab:fb-macro}。这些观察说明需要同时报告问题分布与关系分布，不能将总体微小正差外推为所有频率层级的收益。

\subsection{训练负例数量与排序敏感性}
\label{sec:k-sensitivity}

增加训练负例会扩大每次候选损失的比较范围，但不保证带来更好的排序。我们固定Full-K9获选的超参数、训练种子及优化步数，仅将训练负例数改为$K=3$或$15$；$K=9$模型和评测复用前述结果。三个模型始终在同十个测试选项上评估，训练负例数不改变测试候选集合。

\begin{table}[htbp]
\caption{FB训练负例数的排序敏感性（\%，seed=2026）。同10,000题及原十选项，K9复用；各K保持相同优化步数，非等计算量。}
\label{tab:fb-k-sensitivity}
\centering\small
\begin{tabular}{rrrr}
\toprule
训练负例$K$ & H@1 & MRR & H@3 \\
\midrule
3 & 95.28 & 97.4101 & 99.63 \\
9 & 94.85 & 97.2499 & 99.75 \\
15 & 95.03 & 97.3181 & 99.63 \\
\bottomrule
\end{tabular}
\end{table}


H@1在三个K点间的跨度为0.43个百分点，随K增加没有单调改善。K3的H@1与MRR最高，K9的H@3略高；不同指标对这三个点的排序并不一致。K3与K15相对K9分别净多43题与18题正确，完整纠错与退化计数见附录~\ref{sec:fb-rank-repro}。当前结果说明，在这一单种子固定配置下，更多训练负例未持续转化为更高的排名质量；它不确定全局最优K，也不检验不同K下Full相对配额匹配控制的优势。各组保持相同优化步数而非相同计算量，且未新增直接解码评价，故不作等成本优劣或生成鲁棒性解释。
